% !TeX program = xelatex
% ModelingPaperKit 示例论文 — 独立可编译
\documentclass[UTF8,a4paper,zihao=-4,fontset=none]{ctexart}

\usepackage{../../core/paperkit-base}
\usepackage[zh]{../../core/paperkit-math}
\usepackage[zh]{../../core/paperkit-utils}
% 在 \usepackage{paperkit-utils} 之后、\begin{document} 之前 \input 本文件
\crefformat{figure}{{图}~#2#1#3}
\crefformat{table}{{表}~#2#1#3}
\crefformat{equation}{{式}~#2#1#3}
\Crefformat{figure}{{图}~#2#1#3}
\Crefformat{table}{{表}~#2#1#3}
\Crefformat{equation}{{式}~#2#1#3}
\crefformat{theorem}{{定理}~#2#1#3}
\crefformat{lemma}{{引理}~#2#1#3}
\crefformat{definition}{{定义}~#2#1#3}
\crefformat{corollary}{{推论}~#2#1#3}
\crefformat{assumption}{{假设}~#2#1#3}
\crefformat{remark}{{注}~#2#1#3}
\crefformat{property}{{性质}~#2#1#3}
\crefformat{example}{{例}~#2#1#3}
\Crefformat{theorem}{{定理}~#2#1#3}
\Crefformat{lemma}{{引理}~#2#1#3}
\Crefformat{definition}{{定义}~#2#1#3}

\usepackage{pgfplots}
\usepackage{csvsimple}
\pgfplotsset{compat=1.18}

\setCJKmainfont[AutoFakeBold=2.2]{SimSun}
\setCJKsansfont[AutoFakeBold=2.2]{SimHei}
\setCJKfamilyfont{zhsong}[AutoFakeBold=2.2]{SimSun}
\setCJKfamilyfont{zhhei}[AutoFakeBold=2.2]{SimHei}
\newcommand{\songti}{\CJKfamily{zhsong}}
\newcommand{\heiti}{\CJKfamily{zhhei}}

\graphicspath{{./figures/}{../dummy_data/}}

\title{合成数据预测建模示例\\{\large ModelingPaperKit Walkthrough}}
\author{示例论文（脱敏）}
\date{\today}

\begin{document}
\maketitle

\begin{abstract}
本文演示 ModelingPaperKit 核心宏包用法：基于 \texttt{examples/dummy\_data/} 合成 CSV
完成问题重述、回归建模、时间序列可视化与分类结果展示；涵盖
\texttt{\textbackslash figplot}、\texttt{\textbackslash cref}、定理环境、算法伪代码与三线表。
\end{abstract}

\keywords{数学建模；示例论文；合成数据；交叉引用}

\tableofcontents
\clearpage

% ============================================================
%  一、问题重述
% ============================================================
\section{问题重述}

某企业需根据历史运营数据预测未来指标走势，并在二维特征空间中对客户进行二分类。
赛题提供脱敏后的三类 CSV（\texttt{time\_series.csv}、\texttt{regression.csv}、\texttt{classification.csv}），数据目录为 \mbox{\texttt{examples/dummy\_data/}}。

核心任务：
\begin{enumerate}
  \item 建立可解释的回归模型，刻画特征强度与目标变量关系；
  \item 分析时间序列趋势并给出短期预测思路；
  \item 设计分类边界，使两类样本可分。
\end{enumerate}

\todo{正式论文请替换为真实赛题背景，勿使用本示例措辞。}

% ============================================================
%  二、模型建立
% ============================================================
\section{模型建立}

\subsection{符号与假设}

\begin{definition}[特征向量]\label{def:feature}
记样本特征 $\vb{x}=(x_1,x_2)^\top\in\R^2$，目标 $y\in\R$，设计矩阵
$\mat{X}\in\R^{n\times p}$，参数 $\vb{\beta}\in\R^p$。
\end{definition}

\begin{assumption}[线性可加噪声]
观测满足 $y_i=\vb{x}_i^\top\vb{\beta}+\varepsilon_i$，且
$\varepsilon_i\sim\mathcal{N}(0,\sigma^2)$ 独立同分布。
\end{assumption}

\subsection{回归模型}

最小二乘目标为
\begin{equation}\label{eq:ols}
  \min_{\vb{\beta}}\ \|\mat{X}\vb{\beta}-\vb{y}\|_2^2 .
\end{equation}

\begin{theorem}[正规方程]\label{thm:main}
若 $\mat{X}^\top\mat{X}$ 可逆，则 \eqref{eq:ols} 的解为
$\hat{\vb{\beta}}=(\mat{X}^\top\mat{X})^{-1}\mat{X}^\top\vb{y}$。
\end{theorem}

\begin{lemma}[数值稳定性]\label{lem:stable}
对中心化后的特征，条件数 $\kappa(\mat{X}^\top\mat{X})$ 显著下降，有利于求解稳定\upcite{smith2024synthetic}。
\end{lemma}

整体流程见 \cref{fig:workflow}；\cref{fig:timeseries} 与 \cref{fig:classification} 共同支撑结果分析。
图号区间见 \ref{fig:workflow}--\ref{fig:classification}（多图可用逗号分隔的 \verb|\cref|，见 core 文档）。

\figplot{建模流程示意}{3.6cm}{fig:workflow}{%
  运行 generate\_figures.py 生成}{fig_workflow.pdf}{%
  节点对应数据预处理、特征构建、训练与验证。}

% ============================================================
%  三、求解与结果
% ============================================================
\section{求解与结果}

\subsection{时间序列分析}

图~\cref{fig:timeseries} 展示合成周度指标；亦可直接用 \texttt{pgfplots} 读取 CSV：

\begin{figure}[H]
  \centering
  \begin{tikzpicture}
    \begin{axis}[
      width=0.88\textwidth,
      height=5.2cm,
      xlabel={周次序号},
      ylabel={指标值},
      grid=major,
      legend pos=north west,
      font=\small,
    ]
      \addplot[blue, thick] table[
        x expr=\coordindex,
        y=value,
        col sep=comma,
      ] {../dummy_data/time_series.csv};
      \addlegendentry{观测值}
      \addplot[gray, dashed] table[
        x expr=\coordindex,
        y=trend,
        col sep=comma,
      ] {../dummy_data/time_series.csv};
      \addlegendentry{趋势项}
    \end{axis}
  \end{tikzpicture}
  \caption{由 CSV 直接绘制的时序曲线（与 \cref{fig:timeseries} 数据源一致）}
  \label{fig:ts-pgf}
\end{figure}

\figplot{合成时间序列曲线}{5.8cm}{fig:timeseries}{%
  Matplotlib 自 CSV 导出}{fig_timeseries.pdf}{%
  数据文件：\texttt{time\_series.csv}。}

\subsection{分类结果}

\figplot{二维分类散点图}{6.2cm}{fig:classification}{%
  类别 0/1 散点对比}{fig_classification.pdf}{%
  数据文件：\texttt{classification.csv}。}

\subsection{回归样本表}

表~\cref{tab:reg-sample} 自 CSV 导入前 6 行；完整数据见 \texttt{regression.csv}。

\begin{table}[H]
  \centering
  \caption{回归数据片段（csvsimple 导入）}\label{tab:reg-sample}
  \begin{tabular}{rrrr}
    \topline
    feature & strength & noise & target \\
    \midline
    \csvreader[
      head to column names,
      filter=\ifnum\csvinputline<7,
      late after line=\\,
    ]{../dummy_data/regression.csv}{%
      \feature=\feature, \strength=\strength, \noise=\noise, \target=\target
    }{%
      \feature & \strength & \noise & \target
    }\\
    \bottomline
  \end{tabular}
\end{table}

对比 \cref{fig:ts-pgf}、\cref{fig:timeseries} 与表~\cref{tab:reg-sample} 可见图文混排与智能引用效果。

% ============================================================
%  四、算法设计
% ============================================================
\section{算法设计}

采用梯度下降求解 \eqref{eq:ols}，流程见算法~\cref{alg:gd}。

\begin{algorithm}[H]
  \caption{批量梯度下降（示例）}\label{alg:gd}
  \begin{algorithmic}[1]
    \Require 设计矩阵 $\mat{X}$，观测 $\vb{y}$，学习率 $\eta$，迭代上限 $T$
    \Ensure 参数估计 $\hat{\vb{\beta}}$
    \State 初始化 $\vb{\beta}\gets \vb{0}$
    \For{$t=1$ \textbf{到} $T$}
      \State $\vb{g}\gets \mat{X}^\top(\mat{X}\vb{\beta}-\vb{y})$
      \State $\vb{\beta}\gets \vb{\beta}-\eta\vb{g}$
    \EndFor
    \State \Return $\vb{\beta}$
  \end{algorithmic}
\end{algorithm}

算法~\cref{alg:gd} 与定理~\cref{thm:main} 配合，可复现表~\cref{tab:reg-sample} 所示回归拟合\upcite{zhang2023guide}。

\begin{remark}[实现提示]
生产环境建议改用正规方程或 \texttt{QR} 分解；本示例仅演示 \texttt{algorithm} 环境汉化。
\end{remark}


\section{小结}
如 \cref{thm:main}、\cref{def:feature}、\cref{lem:stable} 与 \cref{fig:workflow}、\cref{tab:reg-sample}、\cref{alg:gd} 所示。
图号区间见 \ref{fig:workflow}--\ref{fig:classification}。
本示例覆盖引擎主要能力。\todo{正式参赛论文请删除 \textbackslash todo 标记。}

\bibliographystyle{plain}
\bibliography{references}

\end{document}
